% ===== §9 Legislation as a flow: contract, morality, self-legislation =====
\section{Legislation as a Flow: Contract, Morality, and the Generative Shaping of the Manifold}
\label{p09:sec:contract}

\noindent
The keystone of \xref{p09:sec:core} named the appropriate poem the non-possessive interpretive cycle that generates real value. But a reader who has followed the semiotics will feel an objection gathering, and it must be met head-on, for unmet it would make the whole distinction between the poem and the contract look too clean, too easily won. \emb{The contract appears to be the purest form of the settling sign.} Its very office is to settle: to fix the terms, to pin the obligations, to declare that the parties are, in this respect, \emph{done}. The marriage registration is a contract; the covenant (\zh{盟书}) is, in a sense, a contract. If contract is settling and settling is pathology (\xref{p09:sec:failures}), then the paper must either reject all contract---which is absurd---or concede that its clean opposition of the poem to the settling sign has missed an entire dimension.

This section confronts the objection by way of a distinction that the author owes to a single corrective insight: that \emb{the contract, too, is a language system---but a language system whose office is to \emph{delimit} generativity rather than to generate content.} The poem is the language system that generates; the contract is the language system that bounds. And just because it bounds, the contract can both stifle generativity and foster it. The section's task is to say which, and why---and, in doing so, to keep the contract \emph{within} the generative framework rather than letting it stand outside as a foundation.

\subsection{The poem and the contract: flow on a manifold, and the shaping of the manifold}
\label{p09:sec:flow-manifold}

Return to the geometric language of \xref{p09:sec:synthesis}. There, everything---the geometric phase, the holonomy, the path---was a \emph{flow on a manifold}. But one question was never put: whence the manifold? Who sets its boundary, its topology, the regions it makes reachable? The answer is the contract.

\begin{claim}[The poem is flow on the manifold; the contract shapes the manifold]
The poem is the language system of generativity: it unfolds a free flow upon the manifold, the phase accruing along the path. The contract is the language system that delimits generativity: it generates no flow of content but shapes the \emph{manifold itself}, constraining the unconstrained free flow into a \emph{structured flow on a manifold} (SFM). The poem answers ``how does the flow unfold upon the manifold''; the contract answers ``what is the shape of the manifold, where lie its boundaries, where its impassable singularities.''
\end{claim}

Let the free generativity be a flow on a manifold $M$. Unconstrained, the flow may reach any region of $M$, including those regions that lead to exhaustion, to extraction, to annihilation---the path on which the rose burns the whole of its resource in a single flowering is also a licit flow on the manifold. The contract is a \emph{structuring} of the manifold: it may excise certain regions (setting impassable singularities---the absolutely forbidden harms), or alter the boundary (setting exit conditions---the reachable edge of the manifold), or reshape the metric and the connectivity (governing toward whom value may flow---who is a licit participant in the cycle of value). Through this structuring the free flow is constrained into a \emph{structured} flow: it remains free within the reshaped manifold $M'$, but can no longer flow into the regions the contract has excised or sealed.\footnote{The geometry can be made precise---the contract as a restriction of the reachable set, the structured flow as a constrained flow on a manifold-with-boundary, the impassable harms as excised singularities---but the rigorous differential-geometric treatment is deferred to the appendix and the Value-Foam study it points to. What matters here is the structural analogy: the contract does not draw the path; it shapes the space within which paths are drawn.}

This at once gives precise sense to ``the contract's path does not matter.'' A good contract does not prescribe which path the flow takes (that would annihilate the holonomy, pressing the poem flat into a contract---the settling failure of \xref{p09:sec:failures}). It shapes only the manifold: it sets the boundary (where the cliff of the exit condition lies), the impassable singularities (which harms are absolutely forbidden), the topology (who is connected to whom, toward whom value may flow). \emb{Within the manifold so shaped, the flow is free, poetic, unprescribed.} The contract and the poem are therefore not opposed: the contract shapes the manifold, and the poem flows upon it freely. The bad contract is the one that prescribes even the steps (the settling failure); the good contract is the one that marks only the manifold's boundary and lets the flow run free within it.

\subsection{The justice of the contract lies in the structure, not the path}
\label{p09:sec:contract-justice}

From this follows the section's sharpest claim, and it relocates the question of justice. We judged the good and vicious cycle, until now, by the character of the \emph{flow}---whether each turn extracted or refluxed, whom it reduced to fuel (\xref{p09:sec:fuel}). The contract directs attention beneath the flow, to the \emph{manifold} on which the flow runs.

\begin{claim}[The justice of the contract is the geometry of the manifold it generates]
The justice of a contract lies in the structure of the manifold it generates, and not in any particular path. A manifold may be unjust in its very geometry---its curvature so distributed that the flow in certain regions must converge toward a centre (structural siphoning), or some position so placed that a being there, however it moves, can only be fuel (the structural pure-fuel position). This is not the viciousness of any single flow but the viciousness of the stage itself: a stage so built that some role, however it plays, is doomed to be extracted from. On an unjust manifold, even the most well-meaning dancer cannot dance a just dance---for the curvature has already bent each of his paths toward exploitation.
\end{claim}

This answers a worry left open in \xref{p09:sec:core}. If the good circle requires that no one be reduced to fuel, but the \emph{structure itself} has placed someone in the position of fuel---a position structurally fixed, regardless of any participant's good will---then the good will of the two (the justice of the flow) cannot save her; the manifold itself must be reshaped (the justice of the structure). This is why the merely moral and the merely poetic do not suffice, and why contract and institution are indispensable: morality and the poem optimize the flow, but only contract and institution can reshape the manifold. It is the reason the appropriate disposition is not always \emph{wu wei}: where the manifold is unjust, to let the flow run free upon it is to consent to the exploitation its geometry dictates. There, the structuring---the legislative---intervention is called for.

\subsection{But the shaping of the manifold is \emph{itself} a flow}
\label{p09:sec:legislation-is-flow}

Here the section must guard against a temptation that an earlier formulation of the argument fell into---the temptation to set ``structure'' against ``flow'' as a separate, deeper stratum, a foundation laid once and standing outside the generativity it grounds. This is false, and the falsity matters, for it is the very thing the paper has criticized throughout: the sanctification of a founding act that stands outside the cycle and declares the loop closed (the settling failure at the level of the ground).

\begin{claim}[The shaping of the manifold is itself a generative flow]
The shaping of the manifold---legislation, covenanting, the laying-down of terms---is not an activity outside the flow but \emph{itself a generative flow}. It unfolds in time; it is generative; it can be poetic or settling, good or vicious, no less than the flow of content upon the manifold. Generativity has two functions: to generate value upon a given structure (the poetic function---flow on the manifold), and to generate or reshape the structure itself (the contractual function---flow \emph{of} the manifold). But the latter is no less a flow. The contract is not a stage standing outside the flow; the contract is \emph{the flow that generates the stage}. It differs only in that its object is the generated structure rather than the content upon the structure---which is why it is more decisive in the generation of structure, as legislation is to adjudication in a society.
\end{claim}

The analogy to legislation and adjudication makes the structure exact, and it is worth drawing out, for it gives the section its spine. \emph{Adjudication} treats a particular case under a given law (a flow upon the manifold)---this is the poetic function: within a given structure, generativity unfolds freely, accruing its phase. \emph{Legislation} generates or amends the law itself (a flow of the manifold)---this is the contractual function: it shapes the structure within which adjudication will unfold. And both are flows; both are generative processes; both unfold in time; both can be good or vicious. Legislation is not a founding myth outside adjudication; it is a sustained, fallible, generative process, exactly as adjudication is. \emb{A society that takes its law for a thing founded once and for all, and a society that takes its law for a living process to be continually re-legislated, differ precisely as the dead contract differs from the living one, as ossification differs from the poem.}

This is the corrective that keeps the contract within the generative framework rather than above it. There is no legislative god standing outside generativity to lay down, once and for all, the manifold upon which all subsequent flow must run. The shaping of the manifold has not escaped the generative criteria; it is subject to them, at its own scale. \emb{Even the making of the stage is a kind of dance.} And this is, once more, a polyphony and a self-reference: even ``the shaping of structure'' has not escaped the criteria of generativity, has found no meta-stratum outside the flow---which is the formal imperative of \xref{p09:sec:conclusions} enacted yet again, at the level of the ground itself.

\subsection{Contract and morality: the hard boundary and the intrinsic curvature}
\label{p09:sec:contract-morality}

Within the contractual flow, a distinction must be drawn that the SFM language renders precise---the distinction between the contract proper (external, enforceable, third-party-compelled) and morality (internalized, self-compelled).

\begin{claim}[The external contract shapes the hard boundary; internalized morality shapes the intrinsic curvature]
The external contract shapes the manifold \emph{externally}: its boundary is maintained by an outside force (the law, the lineage, the sanction of society)---a \emph{hard boundary}, against which the flow that runs is met with external penalty. Internalized morality shapes the manifold \emph{intrinsically}: the boundary has become a part of the subject's own dynamics, needing no external maintenance---an \emph{intrinsic curvature}, by which the subject simply does not flow toward the forbidden region. A relation sustained by pure contract is one whose incentive-compatibility is propped by external penalty---structurally, the temperate version of the counterfeit closure of \xref{p09:sec:antithesis}, sustained by an external injection. A relation whose morality is internalized is one whose curvature is, of itself, positive.
\end{claim}

Hence morality is, in a sense, superior to contract---because it is an endogenous positive curvature rather than an externally imposed boundary-correction. But---and this is the necessary balance---where morality is not yet internalized, or where power is gravely asymmetric, the contract is the indispensable floor. The ideal trajectory is therefore one in which the contract serves as a scaffold (the external hard boundary), morality is gradually internalized (the endogenous curvature generated), and the contract at last \emph{retires} into the background---as, in \xref{p09:sec:repair}, repair gives way again to \emph{wu wei}. \emb{The best contract is the one that makes itself ever less needful of being invoked; the highest achievement of a contract is that it never has to be enforced.}

\subsection{Self-legislation: who may shape the manifold}
\label{p09:sec:self-legislation}

The deepest stratum of the cycle's justice is now reached, and it gathers the whole section. We said that the justice of the structure lies in the geometry of the manifold; but \emph{who shapes the manifold}? If the strong alone shape it, and the weak merely flow upon a manifold given to them, then the very building of the stage is unjust---however symmetric the stage may appear. This carries forward the relational self-legislation of \pinkref{Paper~III}: that the norms of an intimate bond are not imposed from without but are the parties' joint self-legislation---a relational transposition of the Kantian thought that moral law binds because it is autonomously self-imposed \parencite{kant1998}.

\begin{claim}[Just self-legislation: no one excluded from the shaping of the manifold]
A just manifold must be \emph{co-shaped}: every participant who generates value upon the manifold ought to be a co-legislator of the manifold, and not merely a given point upon it. This drives the ``no one is pure fuel'' of \xref{p09:sec:coupled} to its deepest level: not only must no one be reduced to fuel \emph{in the flow}, but no one must be excluded from the shaping \emph{of the manifold}. Self-legislation is the relational co-shaping of the very stage upon which all will generate value. The vicious counterpart is the manifold shaped by the strong alone, presented to the weak as a given---the structural analogue of the demagogue's grammar with but one legitimate derivation (\xref{p09:sec:non-possess}): a stage that looks like a shared agreement while in truth being the strong party's settling, hidden within the illusion of ``our covenant.''
\end{claim}

And self-legislation is, in turn, \emph{itself poetic}---which closes the section back upon the semiotics. The manifold is not shaped once and fixed forever; good self-legislation continually \emph{reshapes} the manifold, as the relation changes, as power shifts, as new participants enter and the manifold must be renegotiated. The dead contract is the manifold ossified after a single shaping (the ossifying failure of \xref{p09:sec:failures}, where ``this is what we agreed at the outset'' becomes a weapon of oppression); the living self-legislation is the manifold continually, jointly re-shaped---\emph{re-legislated like a poem that is never read out}. \emb{So the highest norm, too, is poetic: the limit of the good contract approaches the poem---a living manifold forever open to being jointly re-interpreted and re-legislated.} The justice of a cycle is therefore not two strata, an upper structure over a lower flow, but one generative justice carried through at two scales: in the flow upon the manifold (no one reduced to fuel), and in the flow that shapes the manifold (no one excluded from its joint generation). To co-author the poem is, at this scale, to co-legislate the manifold; and the stage, like the dance, is never finished, and belongs to no one.
